\clearpage
\fancyhead[L]{\small ASTRONOMISCHES JAHRBUCH.}
\hypertarget{native-616-lower}{}
\begin{center}\small Göttingische gelehrte Anzeigen. Stück 80. Seiten $793\;\therefore\;800$. 20. Mai 1819.\\[3pt]\rule{.72\linewidth}{.3pt}\end{center}

\emph{Astronomisches Jahrbuch für das Jahr {\upshape 1821} nebst einer Sammlung der neuesten Abhandlungen, Beobachtungen und Nachrichten aus dem Gebiet der astronomischen Wissenschaften. Berechnet und herausgegeben von} Dr. J. E. \textsc{Bode}, \emph{königlicher Astronom usw.} Berlin 1818. Beim Verfasser und in Kommission bei Ferd. Dümmler. 250 Seiten im Oktavformat, nebst einer Kupfertafel.

\begingroup\brokenpenalty=10000\relax
Da die Einrichtung des astronomischen Kalenders unverändert geblieben ist, brauchen wir nur über die begleitenden Aufsätze zu berichten. Der erste von Prof. \textsc{Degen} in Kopenhagen behandelt das \textsc{Kepler}'sche Problem und entwickelt die Koeffizienten der Reihe für die Mittelpunktsgleichung um zwei Ordnungen weiter als der im vorhergehenden Jahrgang enthaltene Aufsatz von \textsc{Schubert} über dieselbe Aufgabe. Wir verweisen dabei auf das, was wir bei der Besprechung des letzteren angemerkt haben. --- Beobachtete Trabantenverfinsterungen, Sternbedeckungen und Planetenoppositionen von Prof. \textsc{Bürg} an der Wiener Sternwarte. --- Beobachtete Sonnenwenden in den Jahren 1812---1817 an der Turiner Sternwarte von \textsc{Plana}. Auch diese an einem \textsc{Reichenbach}'schen Multiplikationskreis mit stehender Säule, von 18 Zoll Durchmesser, angestellten Beobachtungen ergeben aus den Wintersonnenwenden eine um 4 Bogensekunden geringere Schiefe der Ekliptik als aus den Sommersonnenwenden. Wenn der Herausgeber in einer Anmerkung sagt, dass diese so bemerkenswerte Erscheinung bereits von Astronomen und Physikern untersucht und erklärt worden sei, und auf S. 197 noch bestimmter behauptet, dass sie von \textsc{Bessel} \emph{durch eine verbesserte Refraktionstafel} beseitigt worden sei, so beruht dies auf einem Irrtum. Denn jener Unterschied verschwindet bei Anwendung der \textsc{Bessel}'schen Refraktionstafel keineswegs, sondern fällt eher noch etwas größer aus als bei Anwendung anderer Refraktionstafeln. --- Beobachtete Planetenoppositionen in Kremsmünster von \textsc{Derflinger}. --- Über verschiedene astronomische \SourceLeaf{617}Gegenstände von \textsc{Littrow}. Der Kürze halber beschränken wir uns hier auf einige Bemerkungen zu den Ergebnissen, die Herr \textsc{Littrow} aus den Königsberger Beobachtungen gewonnen hat. Das bereits früher in der Zeitschrift für Astronomie mitgeteilte Verzeichnis der von Herrn \textsc{Littrow} aus den Königsberger Beobachtungen abgeleiteten Polardistanzen von 23 Fixsternen ist hier erneut abgedruckt. Bekanntlich sind diese Polardistanzen alle um mehrere Bogensekunden größer als diejenigen, die \textsc{Piazzi}, \textsc{Oriani} und \textsc{Pond} in großer Übereinstimmung aus ihren Beobachtungen abgeleitet haben. Es bleibt zu hoffen, dass die Ursache dieses auffallenden Unterschieds in Zukunft gefunden wird und man daraus Lehren für die praktische Astronomie ziehen wird. Doch so sehr wir auch die ausgezeichnete Sorgfalt anerkennen, mit der der vortreffliche Königsberger Astronom die Teilungsfehler seines Kreises bestimmt hat, so erscheint uns das Urteil Herrn \textsc{Littrow}s, dass die Königsberger Beobachtungen vor allen anderen von allen konstanten Fehlern ganz freizusprechen sind, zu voreilig und gewagt. Herr \textsc{Littrow} teilt hier ferner die von ihm aus den beiden ersten Jahrgängen der Königsberger Beobachtungen abgeleiteten Rektaszensionsunterschiede der 36 \textsc{Maskelyne}'schen Fundamentalsterne mit, zusammen mit den Gewichten und wahrscheinlichen Fehlern, die diesen Bestimmungen beizulegen sind. Hätten diese Bestimmungen wirklich die Zuverlässigkeit, die ihnen hier zugeschrieben wird, so besteht kein Zweifel, dass sie alles, was wir bisher in dieser Art besaßen, weit übertreffen würden. Doch gerade diese unerhörte Genauigkeit, die man sonst mit den bisherigen Hilfsmitteln nur durch Beobachtungen über die Dauer eines Menschenalters erreichen könnte und die angeblich aus zweijährigen Beobachtungen erreicht worden sein soll, macht eine schärfere Prüfung umso notwendiger, da dieser Gegenstand für die praktische Astronomie von höchster Wichtigkeit ist. Das neueste Heft der Zeitschrift für Astronomie, in dem Herr \textsc{Littrow} sein Verfahren näher angegeben hat, ermöglicht diese Prüfung. Wir sehen daraus, dass Herrn \textsc{Littrow}s Behandlung der Beobachtungen schon für sich genommen dem Geist der Methode der kleinsten Quadrate nicht entspricht und, was noch viel wichtiger ist, dass seine Bestimmung des Gewichts der Ergebnisse diesem Geist in mehreren wesentlichen Punkten ganz entgegensteht und dadurch illusorisch ist. Herr \textsc{Littrow} setzt (Zeitschrift für Astronomie, 6. Band, S. 23) das Gewicht der Bestimmung des Rektaszensionsunterschieds zweier Sterne aus ihren Rektaszensionsunterschieden zu einem dritten, wenn die letzteren Bestimmungen die Gewichte $p$ und $q$ haben, $=\dfrac{4pq}{p+q}$, was falsch ist. Die richtige Formel lautet nur $\dfrac{pq}{p+q}$. Schon dies allein würde die Gewichte seiner Endergebnisse auf fast ein Viertel ihres angeblichen Werts verringern. Noch stärker fällt jedoch der von Herrn \textsc{Littrow} übersehene Umstand ins Gewicht, dass die Verbindung der Teilergebnisse zu einem Endergebnis und die Bestimmung des Gewichts des letzteren nur unter der Voraussetzung gültig sind, dass jene alle ganz unabhängig voneinander sein müssen, was hier keineswegs der Fall ist. Herrn \textsc{Littrow}s große Zahlen werden dadurch so stark herabgesetzt, dass z.\,B. die Bestimmung des Rektaszensionsunterschieds zwischen $\gamma$ im Pegasus und $\alpha$ im Schwan, die angeblich das Gewicht von 1507 einzelnen Vergleichen haben soll, schwerlich viel mehr als das Gewicht von 40 einzelnen Vergleichen behalten dürfte. Auch gegen einige andere Punkte der Rechnung von Herrn \textsc{Littrow} ließen sich Einwände erheben, die wir hier aber der Kürze halber übergehen müssen. --- Opposition der Pallas 1816, und andere astronomische Beobachtungen 1817 von \textsc{David} und \textsc{Bittner} an der Prager Sternwarte. Wir geben hier nur einige Bemerkungen über die mit einem 12-zölligen Multiplikationskreis von \textsc{Reichenbach} beobachteten Zenitdistanzen der Sonne und von Fixsternen. Beobachtungen dieser Art sind auch in mehreren früheren Bänden des Jahrbuchs mitgeteilt worden. Es wäre zu wünschen gewesen, dass diese Beobachtungen mehr in tabellarischer Form zusammengestellt worden wären. Das hätte den Überblick darüber erleichtert, welche Schlüsse sich daraus etwa ziehen lassen. Herr \textsc{David} pflegt die beobachtete Zenitdistanz mit derjenigen zu vergleichen, die er aus der als bekannt angenommenen Polhöhe und der den Sternverzeichnissen oder den Sonnentafeln \SourceLeaf{618}entnommenen Deklination des Gestirns berechnet. Den Unterschied betrachtet er als die beobachtete Refraktion und die Abweichung dieser Refraktion von der Tafelrefraktion als eine Korrektur der letzteren. Dieses Verfahren können wir nicht billigen. Die Polhöhe von Prag hat Herr \textsc{David}, soweit wir wissen, nach der \textsc{Hell}'schen Methode bestimmt, und zwar anhand von Sternen, die im Norden und Süden in beinahe gleicher Höhe kulminieren und deren Deklinationen er den Sternverzeichnissen entnimmt. Das Verfahren ist gut, aber nur für Sternwarten und Instrumente von geringerem Rang. Wer eine so außerordentlich heikle Aufgabe wie die Verbesserung der Refraktionstafeln unternimmt, muss zuvor seine Polhöhe selbstständig und unabhängig von fremden Bestimmungen festsetzen. Was Herr \textsc{David} als Fehler der Refraktionstafeln betrachtet, ist aus dem Zusammenwirken des Fehlers der Polhöhe, der Fehler der Deklinationen, der Fehler des Instruments, der Fehler der Beobachtungen, der Fehler bei den Aberrationsreduktionen und den übrigen Reduktionen und, vielleicht zu einem sehr kleinen Teil, der Fehler der Refraktionstafeln entstanden. Einige wenige isolierte Beobachtungen reichen nicht aus, um diese Fehler voneinander zu trennen. Erforderlich ist eine Reihe von Beobachtungen, die jahrelang mit den besten Instrumenten planmäßig ausgeführt und genutzt wird. Herr \textsc{David} glaubt, dass sich, selbst wenn man die absolute Richtigkeit der Abstände selbst bezweifeln wollte, in den relativen Verhältnissen der beobachteten Zenitdistanzen doch der Einfluss der Refraktionen zeige, die den Temperaturänderungen entsprechend variieren. Diese Änderungen werden jedoch schon bei der Anwendung der Tafeln selbst berücksichtigt. Es ist theoretisch bewiesen, dass, wie auch immer sich die Dichte der Luftschichten ändert, die Refraktion bei so mäßigen Zenitdistanzen nur von der Dichte der untersten Luftschicht abhängt, vorausgesetzt, dass alle Schichten horizontal sind. Dass diese Voraussetzung nicht immer streng erfüllt ist, steht außer Zweifel. Anstelle einer Refraktion, die den Temperaturänderungen folgt, ergibt sich \emph{daraus} jedoch eine unregelmäßige Refraktion, die keinem Kalkül unterworfen werden kann. Ihre Nachteile lassen sich nur durch häufige Beobachtungen und durch Ausschluss derjenigen beheben, bei denen die Gestirne besonders unruhige Bilder gaben. Wir erkennen daher in den erwähnten Unterschieden nur das Spiel der unvermeidlichen Unregelmäßigkeiten der Beobachtungen. Uns ist außerdem aufgefallen, dass Herr \textsc{David} die mit dem Kristallwürfel (Prisma) beobachteten Zenitdistanzen immer um $1''$ vermindert. Wir können diese angebliche Korrektur nicht billigen, da wir in der Natur des Instruments keinen Grund finden, der die Annahme eines für alle Zenitdistanzen konstanten Fehlers rechtfertigt. --- Ferner von derselben Sternwarte Beobachtungen des Jupiter, Uranus, Saturn und der Ceres. --- Komet vom 1. November 1817, der von Dr. \textsc{Olbers} entdeckt, aber nur ein einziges Mal beobachtet wurde. --- Beobachtungen des (ersten) Kometen von 1818 von Dr. \textsc{Olbers}. --- Über einige bemerkenswerte Stellen der Milchstraße von Dr. \textsc{Herschel} (Auszug aus den Philosoph. Transact.). --- Beobachtungen von Planetenoppositionen usw. in Wilna von \textsc{Sniadecki}. --- Beobachtungen und Bestimmungen der Bahn des ersten Kometen von 1818 von \textsc{Enke}. Bemerkenswert an diesem Kometen ist seine Lichtschwäche, die zunahm, als sein Licht (mochte man es als eigenes oder geborgtes betrachten) nach der Rechnung noch hätte zunehmen sollen. Wahrscheinlich erfuhr der Komet während der Dauer seiner Sichtbarkeit bedeutende innere Veränderungen. --- Beobachtete Trabantenverfinsterungen, Sternbedeckungen und Sonnenflecken im Jahr 1817 von \textsc{Hallashka} in Prag. --- Formeln von \textsc{Bessel} zur Berechnung der Position eines Gestirns anhand beobachteter Alignements mit vier Sternen. Sie beruhen im Wesentlichen darauf, zunächst die Schnittpunkte der beiden Großkreise, auf denen die beiden bekannten Sternpaare liegen, mit dem Äquator (oder der Ekliptik) zu bestimmen, obwohl die Lösung eine elegante, rein analytische Form hat. --- Vom selben Astronomen Beobachtungen der Planetenoppositionen, der Sternbedeckungen, der Sonnenwenden und des Polarsterns im Jahr 1817. --- Auszug aus den astronomischen Beobachtungen an der Berliner Sternwarte im selben Jahr. --- Bestimmung der Polhöhe von Krakau. Mit \SourceLeaf{619}einem Repetitionstheodoliten von \textsc{Reichenbach} fand Herr Kreiskommissar \textsc{Lorenz} aus Beobachtungen des Nordsterns $50^\circ\,3'\,57''7$ und aus Sonnenbeobachtungen $50^\circ\,3'\,50''$. Mit dem Mittelwert stimmt die ältere Bestimmung von Herrn \textsc{Sniadecky} sehr genau überein. \textsc{Littrow} hatte sie später um $17''$ vermindert, und zwar, wie aus den neuen Beobachtungen hervorzugehen scheint, zu Unrecht. --- Über die Genauigkeit der astronomischen Beobachtungen und Berechnungen. Dieser Aufsatz soll anschaulich machen, wie viel schwerer es heute ist, in der Astronomie neue Fortschritte zu erzielen als zu der Zeit, da die Anforderungen an den beobachtenden und rechnenden Astronomen so viel geringer waren. Man könnte, sagt der Verfasser, die älteren Astronomen um ihre Lage beneiden. Sie konnten den Lauf der Himmelskörper mit weit weniger Anstrengung beobachten und berechnen, da man noch keine große oder gar höchste Genauigkeit erwarten konnte, und wurden dennoch bei ihren Zeitgenossen und bei der Nachwelt berühmt usw. So wahr das ist, wenn wir die Leistungen der älteren Astronomen von \emph{unserem} Standpunkt aus betrachten, so darf man doch nicht außer Acht lassen, wie unhandlich und unvollkommen ihre Instrumente, wie mühsam ihre Beobachtungsmethoden und wie wenig entwickelt die Rechenverfahren waren. Berücksichtigt man das gebührend, so dürfte beispielsweise \textsc{Tycho} seinen verdienten Ruhm als praktischer Astronom kaum leichter erlangt haben als \textsc{Piazzi}. Zu einzelnen Äußerungen dieses Aufsatzes über die heutige praktische Astronomie ließe sich noch manches einwenden. Nachdem der Verfasser etwa berechnet hat, dass bei einem Kreis von 12 Zoll Radius ein Strich von der Dicke eines Hundertstels einer Linie bereits mehr als 14 Bogensekunden überdeckt, stellt er die Frage „Sollte eine Menschenhand eine Duodezimallinie noch in 100 erkennbare Striche oder Punkte einteilen können?“ Doch eine tatsächliche Teilung der Dezimallinie durch 100 Striche wird gar nicht verlangt. Dass aber die besten Instrumentenbauer die Teilung in die zweckmäßige Anzahl von Teilen so genau ausführen können, dass die Fehler \emph{viel} weniger als ein Hundertstel einer Linie betragen, beweisen die Meisterwerke von \textsc{Reichenbach}, \textsc{Repsold} und \textsc{Troughton} \emph{durch die Tat}. --- Neue Elemente der Junobahn, Beobachtungen dieses Planeten und der Sonne sowie andere astronomische Nachrichten von \textsc{Nicolai}. In diesen Nachrichten wird der Kreis von \textsc{Repsold} an der hiesigen Sternwarte erwähnt. Der Durchmesser dieses Instruments, der vielleicht infolge eines Druckfehlers mit 5 Fuß angegeben ist, ist jedoch auf $3\tfrac{1}{2}$ Pariser Fuß zu berichtigen. --- Beobachtung der Mondfinsternis vom 20. April 1818 in Dresden. --- Rektaszensionen von 35 der wichtigsten Sterne für 1818 nach den letzten Beobachtungen von \textsc{Maskelyne}, aus dem Nautical Almanac für 1820. (Wir bemerken hierzu, dass der Nautical Almanac für 1821 ein noch umfangreicheres Verzeichnis enthält, das auf den Beobachtungen von \textsc{Pond} mit dem neuen Mittagsfernrohr beruht.) Beobachtungen der scheinbaren Helligkeit verschiedener Sterne aus den Jahren 1704---1709 (aus den hinterlassenen Papieren des Berliner Astronomen \textsc{Gottfr. Kirch}). --- Über den Meridiankreis von \textsc{Repsold}, nebst verschiedenen astronomischen Beobachtungen von \textsc{Gauss}. --- Beobachtungen der Sonnenfinsternis vom 5. Mai 1818 von \textsc{Struve}. --- Beobachtung der Opposition der Vesta im Jahr 1818 von \textsc{Enke}. --- Ferner Ephemeriden für Vesta, Juno und Pallas von \textsc{Enke}, \textsc{Posselt}, \textsc{Nicolai} und \textsc{Dirksen}. --- Der Raum erlaubt es nicht, die übrigen kleineren Artikel und Notizen hier einzeln anzuführen.
\par\clearpage\endgroup
